%-------------------------
% Resume in Latex
% Original Author : Sourabh Bajaj
% Adaptation : Hyunggi Chang
% Draft content : Changyeob Lee
% License : MIT
%
% 직군별 이력서(resume_{dfir,malware,cert}_{ko,en}.tex)가 함께 쓰는 서식입니다.
%------------------------

\documentclass[letterpaper,11pt]{article}

\usepackage{latexsym}
\usepackage[empty]{fullpage}
\usepackage{titlesec}
\usepackage{marvosym}
\usepackage[usenames,dvipsnames]{color}
\usepackage{verbatim}
\usepackage{enumitem}
\usepackage[hidelinks]{hyperref}
\usepackage{fancyhdr}
\usepackage[english]{babel}
\usepackage{tabularx}
\usepackage{amsmath}
\usepackage{kotex}
\usepackage{needspace}

\pagestyle{fancy}
\fancyhf{}
\fancyfoot{}
\renewcommand{\headrulewidth}{0pt}
\renewcommand{\footrulewidth}{0pt}

\addtolength{\oddsidemargin}{-0.5in}
\addtolength{\evensidemargin}{-0.5in}
\addtolength{\textwidth}{1in}
\addtolength{\topmargin}{-0.5in}
\addtolength{\textheight}{1.0in}

\urlstyle{same}
\raggedbottom
\raggedright
\setlength{\tabcolsep}{0in}

\titleformat{\section}{
  \vspace{-4pt}\scshape\raggedright\large
}{}{0em}{}[\color{black}\titlerule \vspace{-2pt}]

\newcommand{\resumeItem}[1]{
  \item\small{{#1 \vspace{-2pt}}}
}

% 쪽 끝에서 소제목만 남지 않도록, 남은 공간이 부족하면 다음 쪽에서 시작합니다.
\newcommand{\resumeSubheading}[4]{
  \needspace{3\baselineskip}\vspace{-1pt}\item
    \begin{tabular*}{0.97\textwidth}[t]{l@{\extracolsep{\fill}}r}
      \textbf{#1} & #2 \\
      \textit{\small#3} & \textit{\small #4} \\
    \end{tabular*}\vspace{-5pt}
}

% 설명이 두 줄 이상이어도 아래 항목과 겹치지 않도록 minipage 대신 일반 문단으로 둡니다.
\newcommand{\resumeProject}[2]{
  \vspace{-1pt}\item
    \textbf{#1} \\
    {\small #2\par}\vspace{-3pt}
}

\newcommand{\resumeTalk}[2]{
  \vspace{-1pt}\item
    \textbf{#1} \\
    {\small\itshape #2\par}\vspace{-3pt}
}

\newcommand{\resumeSkill}[2]{
  \item\small{\textbf{#1}: #2}
}

\newcommand{\resumeAward}[3]{
  \item\small{\textbf{#1} / #2 / #3 \vspace{-2pt}}
}

\renewcommand{\labelitemii}{$\circ$}
\newcommand{\resumeSubHeadingListStart}{\begin{itemize}[leftmargin=*]}
\newcommand{\resumeSubHeadingListEnd}{\end{itemize}}
\newcommand{\resumeItemListStart}{\begin{itemize}}
\newcommand{\resumeItemListEnd}{\end{itemize}\vspace{-5pt}}
